%% AMS-LaTeX Created with the Wolfram Language : www.wolfram.com

\documentclass[12pt]{article}

\usepackage[letterpaper,margin=1in]{geometry}
\usepackage{amsmath, amssymb, graphics, setspace}
\usepackage{graphicx} 
\graphicspath{ {./figs/} }
\usepackage{tcolorbox}
\usepackage[framed,numbered,autolinebreaks,useliterate]{mcode}
\newcommand{\mathsym}[1]{{}}
\newcommand{\unicode}[1]{{}}

\usepackage{hyperref}
\hypersetup{
    colorlinks=true,
    linkcolor=blue,
    filecolor=magenta,      
    urlcolor=cyan,
    pdftitle={Overleaf Example},
    pdfpagemode=FullScreen,
}

\def\mathbi#1{\textbf{\em #1}}

\begin{document}
\doublespacing	

\title{CE 401/5501 Homework 09\\
\vspace{2in}
\textbf{Name: \rule{2in}{0.15mm}}}
\author{}
\maketitle

\newcommand*{\I}{\imath} % imaginary unit i

\newpage
\noindent Please review lecture notes \href{https://umsystem.instructure.com/courses/280496/files/33335361?module_item_id=9867784}{Topic 7} and complete the following ML model design problems. 

\section*{Problem 1: Multi-class Classification}

Consider the task of designing a Machine Learning (ML) model to evaluate the health of bridge structures. Based on collected inspection data, the model is expected to classify the health of a bridge into three levels: Excellent Condition, Safe Condition, and Unsafe Condition.

\begin{enumerate}
    \item The ML engineer decides to use a one-hot encoding scheme to represent the three health levels. Assign an exclusive one-hot vector to each of the three conditions. Write down these one-hot vectors.
    
    \item A softmax function as the last layer of the model is implemented. Assume that the model’s decision function outputs three scores: \(z_1 = 10\), \(z_2 = 5\), and \(z_3 = -5\). Calculate the softmax output for these scores. The softmax function for a score \(z_i\) from a vector of scores is given by
    \[
    \text{softmax}(z_i) = \frac{e^{z_i}}{\sum_{j} e^{z_j}},
    \]
    where the sum in the denominator is over all the scores in \(Z\).
    
    \item Based on the softmax output, determine the predicted class label for the bridge's health.
\end{enumerate}

\textbf{Note:} Ensure your solutions include clear calculations and reasoning for each part of the problem.

\newpage

\section*{Problem 2: Analysis of Model Performance}

Given a Machine Learning (ML) model aimed to categorize the condition of bridge structures into three levels: Excellent Condition (level 1), Safe Condition (level 2), and Unsafe Condition (level 3), you are tasked to test the model on a dataset of 100 data points. However, the distribution of the testing dataset is heavily skewed, with 70 data points labeled as Excellent Condition, 25 as Safe Condition, and 5 as Unsafe Condition. 

Additionally, for comparison, a random guessing model is designed, which ignores input data and randomly assigns one of the three labels to each data point.

Two confusion matrices are obtained from evaluating the two models based on the imbalanced testing dataset. Note that for the matrix indices (i.e., Row 1, 2, and 3; Colum 1, 2, and 3), they correspond to the damage levels, level 1, 2, and 3, respettively. 

\begin{enumerate}
    \item The ML model's confusion matrix is given by:
    \[
    \begin{bmatrix}
    60 & 8 & 2 \\
    10 & 14 & 1 \\
    0 & 3 & 2 \\
    \end{bmatrix}
    \]
    
    \item The random model's confusion matrix is given by:
    \[
    \begin{bmatrix}
    24 & 8 & 0 \\
    23 & 9 & 1 \\
    23 & 8 & 4 \\
    \end{bmatrix}
    \]
\end{enumerate}

Please answer the following questions:

\begin{itemize}
    \item [(1)] Verify the validity of the two confusion matrices. Do they correctly reflect the distribution of the testing data?
    \item [(2)] Calculate the Overall Accuracy (OA) from the two confusion matrices. Which model, if only based on the OA measurements, appears better?
    \item [(3)] For the class label 'Excellent Condition', calculate the performance measures of Precision and Recall.
    \item [(4)] For the class label 'Unsafe Condition', calculate the performance measures of Precision and Recall.
    \item [(5)] Can we categorically label the random-guess model as poor? Does it achieve a higher recall for the 'Unsafe Condition' class, and if so, what might be the cause?
\end{itemize}

\end{document}
